\section{Introduction}
\label{sec:introduction}

\subsection{Motivation}

Poker has served as a canonical benchmark for game-theoretic AI since the
early work on Libratus \citep{brown2017libratus} and DeepStack
\citep{moravcik2017deepstack}. Despite significant progress, the practical
relationship between abstraction design and solution quality remains
under-characterized in accessible literature. Most published CFR analyses are
either purely theoretical (convergence guarantees under idealized conditions)
or proprietary (commercial solvers with closed implementations).

Regret addresses this gap. By open-sourcing both the solver and the empirical
analysis framework, we aim to make the following questions answerable without
access to proprietary tools:
\begin{itemize}
  \item How quickly does CFR$^+$ converge as a function of abstraction
        granularity?
  \item How much solution quality do we sacrifice by using coarse card or
        action abstractions?
  \item What is the practical relationship between CFR$^+$ and no-regret
        online learning theory?
\end{itemize}

\subsection{Contributions}

\begin{enumerate}
  \item An open-source CFR$^+$ reference implementation in Rust with
        documented abstractions and a reproducible benchmarking harness.
  \item An empirical convergence analysis measuring exploitability across
        iteration counts, bucket sizes, and bet-size configurations.
  \item An abstraction-error quantification: how abstraction granularity
        affects distance to the true Nash equilibrium.
  \item A theoretical framing connecting CFR$^+$ to FTRL and Blackwell
        approachability, with the regret-flooring step situated within the
        broader online learning literature.
\end{enumerate}
